%% kapa: exam=1 round=19 year=2008 period=2 form=A subject=english subject_ko=영어 seq=28 number=68
%% source: Q-Net 감정평가사 기출문제 (cst00309) artlSeq=1001962 file=extracted/1008/1001962/2008년 제19회 감정평가사 1차시험 기출문제/2교시.hwp
%% conversion: kordoc 4.24.51 -> tools/gen1.py (auto); review=proofread-sonnet; flags=none
\begin{question}{68}
다음 글의 주제로 가장 알맞은 것을 고르시오.
\begin{framedtext}
There are some libraries around you. Now, imagine a library with millions of books――a place where anyone can bring in materials and place them on any shelf or even toss them randomly on the floor. That's something like the way information piles up on the Internet. The Internet, also called the information superhighway, gives you access to millions of pages of text, animation, audio clips, video clips, and more. However, sources of information on the Internet range from the reliable to the suspicious. This fact underscores the need for critical thinking about every aspect of the Internet.
\end{framedtext}
\begin{choices}
  \item Lack of e-books on the Internet
  \item Easy access to information on the Internet
  \item Development of the information superhighway
  \item Need for critical thinking on online information
  \item A variety of materials on the Internet
\end{choices}
\end{question}
\answer{4}
